\subsection{Point orbits and quadratic self-twists}
\label{sec:cartan-image-constraints}

This subsection distinguishes the information carried by a primitive-point
orbit from the stronger existence of a quadratic self-twist of the entire
finite representation.

\begin{proposition}[Maximal point orbits with noncommutative image]
\label{prop:cartan-orbit-data}
Let $p$ be odd, $n\ge2$, $\Lambda_n=\mathbf Z/p^n\mathbf Z$, and let
$D$ be any subgroup of $1+p\Lambda_n$. Put
\[
 U_n(D)=\{g\in\mathrm{GL}_2(\Lambda_n):g\equiv I\pmod p,
                                  \ \det(g)\in D\}.
\]
At each level $1\le j\le n$, the orbit of every primitive vector
$v\in \Lambda_j^2$ is exactly $v+p\Lambda_j^2$, of size $p^{2j-2}$.
Moreover
\[
                  |U_n(D)|=p^{3n-3}|D|,
\]
and $U_n(D)$ is noncommutative when $n\ge3$.
In particular, at $n=4$, the orbit sizes $p^2,p^4,p^6$, even together
with a prescribed determinant image, do not force containment in a
Cartan subgroup.
\end{proposition}

\begin{proof}
The residue of every vector is fixed. To prove transitivity on its
residue fibre, it suffices to use the determinant-one congruence
subgroup. First take $v=e_1$. Any target in this fibre has coordinates
$(a,c)^t$, with $a\equiv1\pmod p$ and $c\equiv0\pmod p$.
Lift $a,c$ to $\Lambda_n$. The matrix
\[
                    \begin{pmatrix}a&0\\c&a^{-1}\end{pmatrix}
\]
has determinant one, is the identity modulo $p$, and carries $e_1$
to the target at level $j$. For arbitrary primitive $v$, lift it
primitively to $\Lambda_n$ and complete it to a basis. Conjugating the
preceding construction by this basis change proves transitivity.
There are $p^{2j-2}$ vectors in the fibre.

The full first congruence subgroup has $p^{4n-4}$ elements: each of
the four entries of $g-I$ is an arbitrary element of $p\Lambda_n$, and
such a matrix is invertible. Its determinant maps onto $1+p\Lambda_n$,
since $\operatorname{diag}(d,1)$ realizes every element $d$ of that
group. Each determinant fibre therefore has $p^{3n-3}$ elements,
which gives the displayed order. Finally, the two elements
$I+pE_{12}$ and $I+pE_{21}$ have determinant one and their products
in the two orders differ by
\[
                         p^2(E_{11}-E_{22}).
\]
This is nonzero modulo $p^n$ when $n\ge3$. Every Cartan arising from
a quadratic commutative algebra is commutative, proving the last
assertion.
\end{proof}

One can also impose the residual markings of the boundary form
without changing this abstract obstruction. The diagonal matrix
$T=\operatorname{diag}(1,-1)$ normalizes $U_n(D)$ and has order two.
The residual representation of
$U_n(D)\rtimes\langle T\rangle$ is split with character ratio $-1$.
Give two copies opposite choices of their distinguished coordinate
line; their identity isomorphism exchanges the distinguished and
complementary markings. Under $U_n(D)$, every lift of either nonzero marked point
still has the stated orbit sizes; the full semidirect product may
have larger orbits because of its tame factor. If $p\equiv1\pmod4$, choose a
root $i\in \Lambda_n$ of $i^2=-1$ and instead take
$T=\operatorname{diag}(i,-i)$. Its determinant is one and its order
is four. The root lifts from $\mathbf F_p$ by successively solving
a linear equation with unit derivative $2i$. Thus even the tame
determinant may be fixed. These are abstract marked representations;
no good formal model or actual unequal-character pair is supplied.

\begin{lemma}[A quadratic self-twist]
\label{lem:quadratic-self-twist-cartan}
Let $p$ be odd, $R=\mathbf Z/p^n\mathbf Z$,
$\rho:\Gamma\to\mathrm{GL}_2(R)$, and let
$\chi:\Gamma\to\{1,-1\}$ be nontrivial. Suppose that an invertible
matrix $J$ satisfies
\[
              J\rho(g)=\chi(g)\rho(g)J\qquad(g\in\Gamma).
\]
Then $\operatorname{tr}(J)=0$ and $J^2=\epsilon I$ for a unit
$\epsilon\in R$. The algebra $R[J]$ is quadratic \'etale, and
\[
                   \rho(\ker\chi)\subset R[J]^\times.
\]
If a subgroup of $\ker\chi$ acts trivially modulo $p$, its image
lies in $1+pR[J]$. This Cartan is nonsplit precisely when the
residue of $\epsilon$ is a nonsquare. A self-twist alone permits
the split case.
\end{lemma}

\begin{proof}
Choose $g$ with $\chi(g)=-1$. Conjugation takes $J$ to $-J$,
so $2\operatorname{tr}(J)=0$. Since $2$ is invertible, the trace
vanishes. Cayley--Hamilton gives
$J^2=-\det(J)I$; take $\epsilon=-\det(J)\in R^\times$.
The reduction of $J$ is not scalar, since an invertible scalar
matrix in odd characteristic has nonzero trace. A nonscalar
$2\times2$ matrix over a field has a cyclic vector: if every
vector spanned a stable line, its action on two basis vectors
and their sum would force it to be scalar. Lift a cyclic vector
$v$ to $R^2$. Then $(v,Jv)$ is a basis.

In this basis $J$ has matrix
$\left(\begin{smallmatrix}0&\epsilon\\1&0\end{smallmatrix}\right)$,
so $R[J]\simeq R[T]/(T^2-\epsilon)$. This is quadratic \'etale,
because $2T$ is a unit in that quotient. If $X$ commutes with $J$
and $Xv=av+bJv$, then $X(Jv)=aJv+b\epsilon v$. Thus
$X=aI+bJ$, proving that the full matrix centralizer is exactly
$R[J]$. Its invertible part is $R[J]^\times$, since the inverse
of a commuting invertible matrix still commutes with $J$.
This proves the asserted containment.

The matrices $I,J$ are linearly independent modulo $p$.
Consequently $aI+bJ\equiv I\pmod p$ forces
$a\equiv1$ and $b\equiv0\pmod p$, giving the principal-unit
containment. If the residue of $\epsilon$ is a square, its two
distinct roots lift and the Chinese remainder theorem gives
$R[J]\simeq R\times R$. If it is a nonsquare, the residue algebra
is $\mathbf F_{p^2}$ and the quadratic \'etale algebra is nonsplit.
Both types are possible: $J=\operatorname{diag}(1,-1)$
anticommutes with the coordinate-swap matrix, giving a split
self-twist example.
\end{proof}

A finite tame extension retains the wild inertia image: its
absolute Galois subgroup contains the subgroup fixing the maximal
tame extension. Thus a tame quadratic self-twist, if actually
identified, places wild inertia in the Cartan of the lemma.
The opposite tangent scalars and exchanged residual lines of
Theorem~\ref{thm:first-boundary-normal-form} do not by themselves
identify the two full deformations as quadratic twists. The
point-orbit restrictions alone are likewise insufficient by
Proposition~\ref{prop:cartan-orbit-data}. Neither elementary result
asserts an actual boundary realization or a boundary character
index.
